% ---------------------------------------------------------------------------
% Phase 2 (Deliverable 3): first-principles EELS/ELNES of the ferroelectric displacement.
% Style-matched to sim_recon_methods.tex / atomfind_methods.tex. Precedes nothing; stands as
% the EELS section. Requires: amsmath, amssymb, graphicx (figs/ generated by an eels figure
% script), booktabs (parameters table). Citations \citep; all keys are defined in
% eels_refs.bib (self-contained). \ref cross-references (sec:sample, fig:tornado, sec:recon,
% sec:atomfind, sec:stem) assume this compiles alongside the sim/atomfind sections.
% Numbers marked <<...>> are filled from the M4/M5 runs when they land.
% ---------------------------------------------------------------------------
\section{First-principles EELS of the ferroelectric displacement}
\label{sec:eels}

\subsection{The model system: a PbTiO$_3$ unit cell and the along-beam question}
\label{sec:eels-sample}

The spectroscopic phase reduces the PbTiO$_3$/SrTiO$_3$ superlattice of Sec.~\ref{sec:sample}
to its ferroelectric building block: a single-formula-unit tetragonal PbTiO$_3$ cell (space
group $P4mm$, no.\ 99), in which the B-site titanium and A-site lead displace from their
centrosymmetric positions along the polar axis $c$, giving the cell an electric dipole. We
quantify that dipole exactly as the imaging pipeline does (Sec.~\ref{sec:atomfind}), by the
titanium offset $\mathbf{u}$ from the charge centre of its oxygen octahedron; at the
experimental lattice ($a=3.90$, $c=4.15$\,\AA, $c/a=1.06$) this is $|\mathbf{u}|\approx
0.33$\,\AA{} along $c$.

The motivation is geometric. Projected STEM imaging and the multislice ptychography of
Sec.~\ref{sec:recon} recover the \emph{in-plane} component of $\mathbf{u}$ --- the polar
displacement projected into the image plane --- but are blind to the component \emph{along the
beam}, which shifts a column along the viewing direction without moving its projected
position. In the polar vortices this hidden component is not incidental: the polarisation
twists with depth into a tube along the beam (Fig.~\ref{fig:tornado}). We therefore ask
whether core-loss EELS, whose near-edge fine structure is sensitive to the local bonding
anisotropy, carries the beam-parallel displacement that projection integrates away. Throughout
we orient the cell to the acquisition of Sec.~\ref{sec:stem} (beam along a pseudocubic
$\langle 100\rangle$), so ``polarisation along the beam'' is the polar axis $c\parallel$ beam
and ``in-plane'' is $c\perp$ beam (Fig.~\ref{fig:eels-model}a).

\subsection{Core-loss EELS and the orientation-dependent near-edge structure}
\label{sec:eels-elnes}

In core-loss EELS the fast electron ionises an atomic core state; the energy-loss near-edge
structure (ELNES), the first ${\sim}30$\,eV above the edge onset, is described by the dynamic
form factor $S(\mathbf{q},E)$ entering the double-differential cross-section
\citep{Egerton2011},
\begin{equation}
  \frac{\partial^2\sigma}{\partial\Omega\,\partial E}
  = \frac{4\gamma^2}{a_0^2\,q^4}\,S(\mathbf{q},E),
  \qquad
  S(\mathbf{q},E) = \sum_f \big\lvert\langle f\lvert e^{\,i\mathbf{q}\cdot\mathbf{r}}\rvert i\rangle\big\rvert^2
  \,\delta\!\big(E - E_f + E_i\big),
  \label{eq:ddscs}
\end{equation}
with $\mathbf{q}$ the momentum transfer, $\gamma$ the relativistic factor, $a_0$ the Bohr
radius, $\lvert i\rangle$ the core state and $\lvert f\rangle$ the unoccupied final states.
For the small $q$ of a near-axis collection the dipole expansion
$e^{\,i\mathbf{q}\cdot\mathbf{r}}\approx 1 + i\,\mathbf{q}\cdot\mathbf{r}$ gives
\begin{equation}
  S(\mathbf{q},E) \;\propto\; q^2 \sum_f
  \big\lvert\langle f\lvert \hat{\mathbf{q}}\cdot\mathbf{r}\rvert i\rangle\big\rvert^2
  \,\delta\!\big(E - E_f + E_i\big),
  \label{eq:dipole}
\end{equation}
i.e.\ the ELNES measures the symmetry-projected unoccupied local density of states resolved
along $\hat{\mathbf{q}}$. In an anisotropic crystal this depends on $\hat{\mathbf{q}}$
relative to the lattice \citep{Nelhiebel1999}; for a small on-axis aperture
$\hat{\mathbf{q}}$ is essentially the beam direction, so the edge shape depends on how the
crystal --- here the polar axis --- is oriented to the beam. This is the observable we exploit.

Two consequences of Eq.~\eqref{eq:dipole} frame the whole study. First, because
$\lvert\langle f\lvert \hat{\mathbf{q}}\cdot\mathbf{r}\rvert i\rangle\rvert^2$ is even in
$\mathbf{q}$, the dipole ELNES is sensitive to the \emph{magnitude} of the along-beam
polarisation but not to its \emph{sign}: reversing $\mathbf{P}$ along the beam is equivalent
to $\mathbf{q}\to-\mathbf{q}$ and leaves Eq.~\eqref{eq:dipole} unchanged. Second, a real
measurement integrates Eq.~\eqref{eq:ddscs} over the convergence ($\alpha$) and collection
($\beta$) apertures; the anisotropic term is weighted by a geometric factor that vanishes at
the \emph{magic angle} $\beta^\star\!\approx 4\,\theta_E$, with the characteristic inelastic
angle $\theta_E = E/(\gamma m_0 v^2)$ \citep{Jouffrey2004}. For the O~K edge at
$300$\,keV, $\theta_E\approx 1.1$\,mrad and $\beta^\star\approx 4.3$\,mrad, far below the
$\alpha=100$\,mrad convergence of the ptychography probe (Sec.~\ref{sec:stem}); the
detectability of the effect in this geometry is quantified in Sec.~\ref{sec:eels-protocol}.

\subsection{First-principles ELNES with core-hole DFT}
\label{sec:eels-dft}

We evaluate Eq.~\eqref{eq:dipole} from first principles with the plane-wave pseudopotential
code CASTEP \citep{Clark2005} using the core-hole method \citep{Gao2008}. The ionised atom is
given a distinct species label and an on-the-fly-generated pseudopotential whose core shell
carries one fewer electron --- a full core hole (occupancy $1s^1$ for the O~K edge, $2p^5$ for
Ti~L$_{2,3}$, $3d^9$ for Pb~M), with a neutralising background (cell charge $+1$) --- so that
the conduction states are computed in the presence of the hole. A $2\times2\times2$ supercell
dilutes the core hole from its periodic images. The core-to-conduction matrix elements written
by CASTEP are then contracted by OptaDOS \citep{Morris2014,Nicholls2012} into the core-loss
spectrum, with the momentum direction $\hat{\mathbf{q}}$ set explicitly and life-time and
instrument (Gaussian/Lorentzian) broadening applied. The quantity central to this work is the
\emph{ELNES dichroism},
\begin{equation}
  \Delta(E) \;=\; S(\mathbf{q}\!\parallel\! c,\,E)\;-\;S(\mathbf{q}\!\perp\! c,\,E),
  \label{eq:dichroism}
\end{equation}
the difference between probing along and across the polar axis; by the reciprocity of
Sec.~\ref{sec:eels-sample} this is exactly the change in the measured edge when the
polarisation is rotated from along the beam to in-plane. We take the O~K edge as the primary
observable --- it is well described at the single-particle level and is an established probe of
the ferroelectric off-centering \citep{TorresPardo2011,Bugnet2016} --- and treat Ti~L$_{2,3}$
for anisotropy trends only, since its $2p$ spin--orbit multiplet structure is not captured by
single-particle DFT \citep{Mizoguchi2009}.

\subsection{Computational protocol and validation}
\label{sec:eels-protocol}

The calculation is built as a validated staircase, each step gated before the next; the code,
structures and per-run inputs are in the project repository (Sec.~\ref{sec:eels-repro}).

\paragraph{Structure and convergence.}
Cells are generated and checked against the experimental crystallography (space group,
$c/a$, $|\mathbf{u}|$) before any DFT. The plane-wave cutoff and $k$-point grid are converged
to $<1$\,meV/atom, and the setup is validated by the ferroelectric \emph{double well}: single-
point energies across a displacement series $\mathbf{u}=s\,\mathbf{u}_0$, $s\in[0,1]$, at fixed
tetragonal strain place the fully polar structure $203$\,meV per formula unit below the
centrosymmetric reference, monotonically, with the minimum at the experimental displacement
($s\!=\!1$) --- confirming that CASTEP/PBE reproduces the ferroelectric instability
(Fig.~\ref{fig:eels-model}b).

\paragraph{Method benchmark.}
The core-hole recipe is fixed by reproducing a known edge: the rutile TiO$_2$ O~K ELNES comes
out with its textbook shape --- the t$_{2g}$/e$_g$ near-edge peaks split by $2.7$\,eV with
$e_g>t_{2g}$, followed by the broad O~$2p$--Ti~$4sp$ band --- matching experiment
(Fig.~\ref{fig:eels-benchmark}) \citep{Mizoguchi2009}. This locks the core-hole occupancy,
the charge treatment and the OptaDOS parameters used for PbTiO$_3$.

\paragraph{Null test and dichroism.}
As an artifact check, the dichroism of the \emph{cubic} cell computed at the isotropic Ti site
must vanish; any residual is numerical. The polar signal is then the dichroism
$\Delta(E)$ of the tetragonal cell, Eq.~\eqref{eq:dichroism}, for the O~K (and Ti~L) edges
(Fig.~\ref{fig:eels-dichroism}a). Rotational invariance is verified by the equality of the
$c\perp$ beam spectrum of the $c\!\parallel\! z$ cell and the $c\!\parallel\!$ beam spectrum
of the $c\!\parallel\! x$ cell. \emph{The along-beam polarisation produces a
<<XX>>\,\% near-edge dichroism, concentrated in the <<pre-peak / e$_g$>> region.}

\paragraph{Displacement dependence (the objective).}
The core deliverable --- how the atomic displacement shapes the spectrum --- is the evolution
of the edge and its dichroism across the series $s\in[0,1]$ (Fig.~\ref{fig:eels-dichroism}b).
The intercept at $s\!=\!0$ isolates the non-polar crystal-field/strain anisotropy of the
oxygen backbone; the slope in $s$ is the ferroelectric contribution, the first-principles
counterpart of the empirical ``broadening grows with Ti off-centering'' of
\citet{Bugnet2016}. \emph{We find the dichroism metric to be <<monotonic / near-linear>> in
$|\mathbf{u}|$, giving the map from a measured anisotropy back to the along-beam displacement.}

\paragraph{Detectability in the acquisition geometry.}
Finally the intrinsic dichroism is averaged over the $300$\,keV collection geometry: the
measurable anisotropy is the surviving fraction after integrating Eq.~\eqref{eq:ddscs} over the
aperture, which crosses zero at the magic angle $\beta^\star\approx4.3$\,mrad
(Fig.~\ref{fig:eels-dichroism}c). The $100$\,mrad convergence of the ptychography probe lies far
past this, averaging the anisotropy away; recovering the along-beam signal in practice requires
a small ($\lesssim2$\,mrad) EELS collection aperture rather than the ptychography optics. This
sets the conditions under which the stretch-goal fusion of EELS with the depth-resolved
ptychography of Sec.~\ref{sec:recon} could recover the full three-dimensional polarisation.

\begin{figure}[t]
  \centering
  \includegraphics[width=\linewidth]{figs/eels_model.pdf}
  \caption{Model system and DFT validation. (a) The $P4mm$ PbTiO$_3$ unit cell (Pb, Ti, O),
    with the Ti off-centering $\mathbf{u}$ along the polar axis $c$; the two orientations used
    are $c\parallel$ beam (``P along beam'') and $c\perp$ beam, matching the acquisition of
    Sec.~\ref{sec:stem}. (b) Ferroelectric double well: total energy versus the scaled
    displacement $s$ at fixed tetragonal strain, the polar structure ($s\!=\!1$) lying
    $203$\,meV/f.u.\ below the centrosymmetric reference with its minimum at the experimental
    displacement. [PLACEHOLDER --- generated by the eels figure script.]}
  \label{fig:eels-model}
\end{figure}

\begin{figure}[t]
  \centering
  \includegraphics[width=0.68\linewidth]{figs/eels_benchmark.pdf}
  \caption{Method benchmark: calculated rutile TiO$_2$ O~K ELNES (core-hole CASTEP + OptaDOS)
    against the reference edge, reproducing the t$_{2g}$/e$_g$ near-edge peaks (split
    $2.7$\,eV, $e_g>t_{2g}$) and the broad O~$2p$--Ti~$4sp$ band. This fixes the core-hole
    recipe used for PbTiO$_3$. [PLACEHOLDER.]}
  \label{fig:eels-benchmark}
\end{figure}

\begin{figure}[t]
  \centering
  \includegraphics[width=\linewidth]{figs/eels_dichroism.pdf}
  \caption{Along-beam polarisation in the O~K ELNES. (a) Edge computed with the momentum
    transfer along ($\mathbf{q}\!\parallel\!c$, P along beam) and across ($\mathbf{q}\!\perp\!c$)
    the polar axis, and their difference, the dichroism $\Delta(E)$ (Eq.~\ref{eq:dichroism}).
    (b) The dichroism metric versus the Ti off-centering $|\mathbf{u}|$ across the displacement
    series, separating the non-polar backbone anisotropy (intercept) from the ferroelectric
    contribution (slope). (c) Fraction of the intrinsic dichroism surviving collection at
    semi-angle $\beta$ (300\,keV), crossing zero at the magic angle
    $\beta^\star\!\approx\!4.3$\,mrad; the $100$\,mrad ptychography convergence averages it out.
    [PLACEHOLDER.]}
  \label{fig:eels-dichroism}
\end{figure}

\subsection{Reproducibility}
\label{sec:eels-repro}

All parameters are collected in Table~\ref{tab:eels-params}. Structure generation and
validation (\texttt{build\_cells.py}), the CASTEP core-hole and OptaDOS jobs
(\texttt{run\_coreloss.slurm}), and the spectral analysis (\texttt{analyze\_elnes.py}) are
scripted end to end; a unit-test suite guards the structure and parsing code. Cells,
templates and per-run inputs are in the repository, so a spectrum is reproduced by
regenerating the cell and re-running the driver with the tabulated values.

\begin{table}[t]
  \centering
  \small
  \begin{tabular}{@{}lll@{}}
    \toprule
    Stage & Parameter & Value \\
    \midrule
    Structure & material / space group & PbTiO$_3$, $P4mm$ (polar), $Pm\bar{3}m$/$P4/mmm$ (ref.) \\
              & lattice; Ti off-centre  & $a{=}3.90$, $c{=}4.15$\,\AA{} ($c/a{=}1.06$); $|\mathbf{u}|{\approx}0.33$\,\AA \\
              & core-hole supercell     & $2\times2\times2$ (40 atoms) \\
              & benchmark               & rutile TiO$_2$ ($P4_2/mnm$), O~K edge \\
    \addlinespace
    DFT       & code / functional       & CASTEP \citep{Clark2005} / PBE \\
              & cutoff; $k$-grid         & <<X>>\,eV; <<n$^3$>> (converged $<1$\,meV/atom) \\
              & core hole               & full (occ.\ $1s^1$/$2p^5$/$3d^9$), cell charge $+1$ \\
              & empty bands             & \texttt{nextra\_bands}\,$=100$ \\
    \addlinespace
    Spectrum  & code / task             & OptaDOS \citep{Morris2014} / core-loss, absorption \\
              & momentum directions     & $\hat{\mathbf{q}}\!\parallel\!c$ and $\hat{\mathbf{q}}\!\perp\!c$ (\texttt{core\_qdir}) \\
              & broadening              & adaptive + LAI (Gauss $1.0$, Lorentz $0.5$\,eV) \\
              & edges                   & O~K (primary); Ti~L$_{2,3}$, Pb~M (secondary) \\
    \addlinespace
    Geometry  & beam energy             & $300$\,keV ($\theta_E^{\mathrm{OK}}\!\approx\!1.1$\,mrad) \\
              & magic angle             & $\beta^\star\!\approx\!4.3$\,mrad \\
              & orientation             & beam $\parallel\langle100\rangle$; $c\parallel$ vs $\perp$ beam \\
    \bottomrule
  \end{tabular}
  \caption{First-principles EELS parameters.}
  \label{tab:eels-params}
\end{table}
% ---------------------------------------------------------------------------
